% gml-layout.tex - the TeX side of gml's layout mode (gml -layout CMD,
% lib/layout.ml).  Input after gml.tex, by the driver file that gml
% writes:
%
%   \def\gmlmode{hints} \input gml \input gml-layout
%   \def\gmllayoutdir{/tmp/dir/}  ... \gmlbegindocument%
%   \gmlloop
%
% TeX must run in scroll mode (or errorstop mode): \read-1 from the
% terminal is fatal in batch and nonstop mode.  \gmlloop writes
% "GML:READY" on the terminal and reads one line, which gml writes:
%
%   \gmlchunk{k}\relax\gmlquery{pn}   input chunk k (the document's TeX
%                                     source so far, whole lines), then
%                                     answer "GML:ANS n" (pn or lc)
%   \gmlchunk{k}\relax\gmlfinish      the last chunk; end the document
%
% \relax ends any number or keyword TeX is scanning at the end of the
% chunk, as the next line of the document would.  The line is read with
% \endlinechar=-1, so no end-of-line token reaches the document (which
% may be in format-off mode, where ^^M is active).  TeX's own terminal
% output is ignored by gml except error lines ("!").
%
%   pn  the page number, \pageno (its absolute value: plain TeX counts
%       roman pages negative).  DCF-VERIFY: &$PN with .pn roman.
%   lc  the lines left on the page (&$LC, "lines left in the column",
%       SX26-3723): (\pagegoal-\pagetotal)/\baselineskip, rounded down;
%       \vsize/\baselineskip on an empty page (\pagegoal=\maxdimen).
%       In outer vertical mode the page builder is run first (\par, which
%       does nothing else there, unless a \parshape, \hangindent or
%       \looseness is set for the next paragraph: then it is not run).  In
%       horizontal mode the paragraph in progress is not counted: forcing
%       \par would change the output.  DCF-VERIFY: DCF counts output lines
%       already formatted, including the partial line in progress.
\ifx\gmlloop\undefined\else\endinput\fi
\chardef\gmllayoutatcode=\catcode`\@
\catcode`\@=11

\newcount\gml@lcnt \newdimen\gml@ldim \newdimen\gml@lbs
\ifx\gmllayoutdir\undefined\def\gmllayoutdir{}\fi

\def\gmlsend#1{\immediate\write16{GML:#1}}
\def\gmlloop{\gmlsend{READY}\gml@lread\gml@lcmd\gmlloop}
\def\gml@lread{\scrollmode{\endlinechar\m@ne\global\read\m@ne to\gml@lcmd}}
\def\gmlchunk#1{\input \gmllayoutdir c#1.tex }

\def\gmlquery#1{\expandafter\ifx\csname gml@q@#1\endcsname\relax
    \gmlsend{ANS ?}\else\csname gml@q@#1\endcsname\fi}
\def\gml@q@pn{\gml@lcnt\pageno \ifnum\gml@lcnt<\z@\gml@lcnt-\gml@lcnt\fi
  \gmlsend{ANS \the\gml@lcnt}}
\def\gml@q@lc{\gml@lbuild
  \ifdim\pagegoal=\maxdimen\gml@ldim\vsize
  \else\gml@ldim\pagegoal\advance\gml@ldim-\pagetotal\fi
  \gml@lbs\baselineskip\ifdim\gml@lbs>\z@\else\gml@lbs12\p@\fi
  \gml@lcnt\gml@ldim\divide\gml@lcnt\gml@lbs
  \ifnum\gml@lcnt<\z@\gml@lcnt\z@\fi
  \gmlsend{ANS \the\gml@lcnt}}
\def\gml@lbuild{\ifvmode\ifinner\else\ifnum\parshape=\z@\ifdim\hangindent=\z@
  \ifnum\looseness=\z@\ifnum\hangafter=\@ne\endgraf\fi\fi\fi\fi\fi\fi}

% An anchor (gml:tcmark: a table of contents entry) also tells gml its
% page when the page is shipped out: "GML:PAGE id folio" on the
% terminal (a \write that is not \immediate is done at \shipout), so
% that a .tc knows the pages of the entries before it without asking in
% the middle of a paragraph.
\let\gml@lanchor\gmlanchor
\def\gmlanchor#1{\edef\gml@tmpa{\write16{GML:PAGE #1 \noexpand\gml@folio}}%
  \gml@tmpa\gml@lanchor{#1}}

% (\bye is \outer in plain TeX)
\def\gmlfinish{\gmlenddocument\csname bye\endcsname}

% A line scope written before its continuation was known (a query made
% gml send it early): more=yes after all, given before \gmlendlines.
\def\gmllinesmore{\edef\gml@thislattrs{\gml@thislattrs,more=yes}}

\catcode`\@=\gmllayoutatcode
